\section{Discussion}
% GOAL: 400-600 words. Interpretation, limitations, future directions.

\subsection{Why Serial Benchmarks Matter}

Our results demonstrate that serial dependency introduces an evaluation axis invisible to static benchmarks.
When tasks are independent, an agent's failure at one task does not affect others, and the benchmark can measure each capability in isolation.
When tasks are serially dependent, failure propagates, and the benchmark must distinguish between ``cannot reach'' and ``cannot solve''---a distinction that resurrection enables.
This is not merely a technical detail but a fundamental shift in what benchmarks measure: from \emph{can the agent solve task~$k$?} to \emph{can the agent progress along the dependency chain, and if not, where and why?}

\subsection{Resurrection as Methodology, Not Assistance}

A natural concern is that resurrection ``helps'' the agent by providing correct intermediate artifacts, thereby inflating scores.
We emphasize that resurrection is an \emph{evaluation protocol}, not a scoring adjustment.
The no-resurrection setting (Setting~2) provides the ``strict'' measurement of zero-shot pipeline capability.
The resurrection setting (Setting~1) provides the ``diagnostic'' measurement of node-level capability given correct prerequisites.
Both are reported independently; neither replaces the other.
The value of resurrection lies in the \emph{comparison} between the two settings, which reveals how much downstream signal is lost to cascading failure.

\textbf{Potential biases.}
We acknowledge three potential biases introduced by resurrection:
(1)~\emph{State discontinuity}: the golden artifact may differ in style or structure from the agent's own output, creating an unnatural transition that could affect downstream performance;
(2)~\emph{Selection bias}: only failed tasks are resurrected, creating a non-random sample of downstream measurements;
(3)~\emph{Moral hazard}: if agents were aware of the resurrection mechanism, they might reduce early-stage effort.
We mitigate bias~(3) by ensuring that agents are \emph{not informed} about resurrection---the mechanism is triggered externally by the evaluation orchestrator, not by the agent itself.
Biases~(1) and~(2) are inherent to the methodology and cannot be fully eliminated; we report them as known limitations and note that the \emph{direction} of bias~(1) is ambiguous---a golden artifact may actually \emph{help} downstream performance (by providing cleaner input) or \emph{hinder} it (by creating unfamiliar state), and our data shows both effects across different models.

\subsection{Limitations}

\textbf{Domain specificity.}
\evobench{} is grounded in compiler construction, a domain with particularly clean intermediate artifact specifications.
Other serial workflows---ML pipelines, DevOps chains, data engineering flows---may have less well-defined intermediate states, making resurrection harder to implement.
We do not claim that \evobench{} generalizes to all serial settings, but argue that the \emph{evaluation methodology} (serial dependency + resurrection + process metrics) is transferable.

\textbf{Task count and statistical power.}
With six tasks, \evobench{} provides limited statistical power for per-task analysis compared to benchmarks with hundreds or thousands of instances.
However, the serial structure means that each task carries \emph{qualitatively different} information (different dependency depth, different failure propagation pattern), which compensates for the small task count with structural diversity.
Additionally, each model configuration is evaluated only once per pipeline run; we do not report variance across multiple seeds.
The primary source of randomness in our evaluation is the LLM's stochastic generation, not the benchmark itself (scoring scripts are deterministic).
We acknowledge that single-run evaluation limits variance estimation and recommend multi-seed evaluation as future work.

\textbf{Contamination risk.}
As with any benchmark based on public code, there is a risk that models have seen the \yatcc{} codebase during training.
\yatcc{} is an open-source course project hosted on GitHub, and frontier models may have been trained on its contents.
However, \evobench{}'s key challenge is not reproducing known solutions but maintaining \emph{serial continuity}---producing correct intermediate artifacts that integrate across the pipeline.
Memorizing individual task solutions does not guarantee correct cross-task artifact flow.
Furthermore, the resurrection setting specifically tests whether agents can \emph{continue} from a given state, which requires integration capability rather than recall.
We note that SWE-bench~\citep{jimenez2024swebench} faces similar contamination concerns and has addressed them through task deduplication and temporal splitting; future versions of \evobench{} could adopt analogous safeguards.

\subsection{Future Directions}

\begin{itemize}[nosep]
\item \emph{Multi-domain serial benchmarks}: extending the serial evaluation methodology to ML engineering, DevOps, and data science pipelines.
\item \emph{Evolutionary gain measurement}: comparing agent performance when inheriting self-produced vs.\ golden intermediate artifacts, quantifying how well agents ``understand their own code.''
\item \emph{Dynamic resurrection}: allowing agents to request resurrection proactively rather than having it triggered automatically, studying self-diagnosis capability.
\item \emph{Expanded model and scaffold coverage}: evaluating additional frontier models and agent frameworks as they become available.
\end{itemize}